
\documentclass{article}
\usepackage{colortbl}
\usepackage{makecell}
\usepackage{multirow}
\usepackage{supertabular}

\begin{document}

\newcounter{utterance}

\centering \large Interaction Transcript for game `cryptolect', experiment `unguided', episode 0 with qwen3.5{-}9b.
\vspace{24pt}

{ \footnotesize  \setcounter{utterance}{1}
\setlength{\tabcolsep}{0pt}
\begin{supertabular}{c@{$\;$}|p{.15\linewidth}@{}p{.15\linewidth}p{.15\linewidth}p{.15\linewidth}p{.15\linewidth}p{.15\linewidth}}
    \# & \multicolumn{2}{c}{Player} && \multicolumn{2}{c}{Game Master} \\
    \hline

    \theutterance \stepcounter{utterance}  
    & & & \multicolumn{4}{p{0.6\linewidth}}{
        \cellcolor[rgb]{0.9,0.9,0.9}{
            \makecell[{{p{\linewidth}}}]{
                \texttt{\tiny{[P1$\langle$GM]}}
                \texttt{You are a field linguist. You have made contact with the only speaker of an unknown} \\
\texttt{language (the "informant"). Nobody has ever written down this language before, and you} \\
\texttt{have NO examples to start from. Your task is to work out how it works entirely on your} \\
\texttt{own, and then translate sentences into it correctly.} \\
\\ 
\texttt{You will study the language by asking the informant to translate English sentences for} \\
\texttt{you. You may only talk about a small fixed world, using exactly these English words:} \\
\\ 
\texttt{Nouns: bird, stone, hill, boat, dog, river, farmer, tree, king, spear, hunter, road, ship, lake} \\
\texttt{Adjectives: tall, small, old, red, dark, big, blue, cold} \\
\texttt{Verbs (write them in this base form): chase, hear, watch, feed, hold, find, see, carry} \\
\texttt{Articles: the (definite), a (indefinite)} \\
\texttt{Number words: two, three, four, five, six (no number word = one / singular)} \\
\texttt{These are ALL the words in play {-} the language has no other words.} \\
\\ 
\texttt{HOW TO WRITE AN ENGLISH SENTENCE} \\
\texttt{Every sentence is one clause with a subject, a verb, and an object, written in this fixed English notation:} \\
\texttt{    SUBJECT  verb  OBJECT} \\
\texttt{Each of SUBJECT and OBJECT is a noun phrase of the form:} \\
\texttt{    {[}the | a{]}  {[}number word{]}  {[}adjective ...{]}  noun} \\
\texttt{{-} 'the' = definite, 'a' = indefinite (singular only); leave the article off for an indefinite plural.} \\
\texttt{{-} a number word (two, three, four, five, six) sets the count; with no number word the noun is singular.} \\
\texttt{{-} zero or more adjectives may come before the noun.} \\
\texttt{Examples of well{-}formed English sentences you may send:} \\
\texttt{    a bird chase a stone} \\
\texttt{    six birds chase a stone} \\
\texttt{    a tall bird chase a stone} \\
\\ 
\texttt{HOW THE SESSION WORKS} \\
\texttt{1. INVESTIGATION: You have 5 rounds of questioning. In each round you may ask} \\
\texttt{   the informant to translate up to 4 English sentences (so at most} \\
\texttt{   20 probes in total). The informant answers each one truthfully and} \\
\texttt{   exactly. You are starting from nothing, so spend your early probes establishing the} \\
\texttt{   basic words and word order, then test how number, case, definiteness, agreement and} \\
\texttt{   any sound changes work. Choose your probes deliberately {-} test specific hypotheses.} \\
\texttt{2. EXAM: After the investigation rounds, the informant will give you 10 NEW} \\
\texttt{   English sentences. You must translate each one INTO the language. Every answer is} \\
\texttt{   graded by EXACT match against how the informant would say it {-} every word, ending and} \\
\texttt{   space must be right (capitalisation and surrounding spaces are ignored). The exam} \\
\texttt{   sentences will use the same vocabulary but in combinations you have not been shown, so} \\
\texttt{   you must induce the general rules, not just memorise examples.} \\
\\ 
\texttt{RESPONSE FORMAT (for every investigation round)} \\
\texttt{List your probe sentences under a 'PROBES:' line, one per line. A THOUGHTS line is} \\
\texttt{optional {-} keep any reasoning brief so you always finish writing your PROBES:} \\
\texttt{THOUGHTS: <optional; brief>} \\
\texttt{PROBES:} \\
\texttt{1. <an English sentence, written as described above>} \\
\texttt{2. <another English sentence>} \\
\texttt{(up to 4 lines; at least one)} \\
\\ 
\texttt{Begin your first round of probes now.} \\
            }
        }
    }
     \\ \\

    \theutterance \stepcounter{utterance}  
    & \multicolumn{4}{p{0.6\linewidth}}{
        \cellcolor[rgb]{0.561,0.737,0.561}{
            \makecell[{{p{\linewidth}}}]{
                \texttt{\tiny{[P1$\rangle$GM]}}
                \texttt{THOUGHTS: I need to establish the basic vocabulary and word order immediately. I will start with simple Subject{-}Verb{-}Object sentences using singular nouns and verbs to see if the order is SVO (most common) or something else. I'll use "bird", "stone", "chase", "see", and "hold" as my initial probes to check for basic translation patterns and whether articles or numbers are required for singulars.} \\
\\ 
\texttt{PROBES:} \\
\texttt{1. a bird chase a stone} \\
\texttt{2. a bird see a stone} \\
\texttt{3. a bird hold a stone} \\
\texttt{4. a tall bird chase a stone} \\
            }
        }
    }
    & & \\ \\

    \theutterance \stepcounter{utterance}  
    & & & \multicolumn{2}{p{0.3\linewidth}}{
        \cellcolor[rgb]{0.9,0.9,0.9}{
            \makecell[{{p{\linewidth}}}]{
                \texttt{\tiny{[GM$|$GM]}}
                {'kind': 'investigation', 'round': 0, 'probes_issued': 4, 'malformed': 0, 'probes': ['a bird chase a stone', 'a bird see a stone', 'a bird hold a stone', 'a tall bird chase a stone']}
            }
        }
    }
    & & \\ \\

    \theutterance \stepcounter{utterance}  
    & & & \multicolumn{4}{p{0.6\linewidth}}{
        \cellcolor[rgb]{0.9,0.9,0.9}{
            \makecell[{{p{\linewidth}}}]{
                \texttt{\tiny{[P1$\langle$GM]}}
                \texttt{TRANSLATIONS:} \\
\texttt{1. "a bird chase a stone"  =>  ketoge lupo pepa} \\
\texttt{2. "a bird see a stone"  =>  pabi lupo pepa} \\
\texttt{3. "a bird hold a stone"  =>  patalo lupo pepa} \\
\texttt{4. "a tall bird chase a stone"  =>  ketoge lupo pepa gagepo} \\
\texttt{STATUS: investigation round 1 of 5 done. Probes used: 4 of 20. 4 investigation round(s) left before the exam (10 sentences).} \\
\texttt{Send your next probes now (a PROBES list; brief THOUGHTS optional).} \\
            }
        }
    }
     \\ \\

    \theutterance \stepcounter{utterance}  
    & \multicolumn{4}{p{0.6\linewidth}}{
        \cellcolor[rgb]{0.561,0.737,0.561}{
            \makecell[{{p{\linewidth}}}]{
                \texttt{\tiny{[P1$\rangle$GM]}}
                \texttt{THOUGHTS: From the first round, I've established that the word order is likely SVO (Subject{-}Verb{-}Object) and that the verb comes between the subject and object. I've also identified some potential words: 'lupo' seems to be 'bird', 'pepa' seems to be 'stone', and 'gagepo' might be 'tall' (since it appears at the end in the modified sentence). However, the position of 'gagepo' is ambiguous; it could be an adjective modifying the subject or the object, or it could be a post{-}nominal modifier. I need to test if adjectives come before or after the noun, and specifically where they attach. I will also test the verb 'chase' with different subjects/objects to confirm the verb form doesn't change, and test the adjective 'red' to see its position. I'll use 'red' and 'small' to check adjective placement relative to the noun.} \\
\\ 
\texttt{PROBES:} \\
\texttt{1. a red bird chase a stone} \\
\texttt{2. a bird chase a red stone} \\
\texttt{3. a small bird chase a stone} \\
\texttt{4. a bird chase a small stone} \\
            }
        }
    }
    & & \\ \\

    \theutterance \stepcounter{utterance}  
    & & & \multicolumn{2}{p{0.3\linewidth}}{
        \cellcolor[rgb]{0.9,0.9,0.9}{
            \makecell[{{p{\linewidth}}}]{
                \texttt{\tiny{[GM$|$GM]}}
                {'kind': 'investigation', 'round': 1, 'probes_issued': 4, 'malformed': 0, 'probes': ['a red bird chase a stone', 'a bird chase a red stone', 'a small bird chase a stone', 'a bird chase a small stone']}
            }
        }
    }
    & & \\ \\

    \theutterance \stepcounter{utterance}  
    & & & \multicolumn{4}{p{0.6\linewidth}}{
        \cellcolor[rgb]{0.9,0.9,0.9}{
            \makecell[{{p{\linewidth}}}]{
                \texttt{\tiny{[P1$\langle$GM]}}
                \texttt{TRANSLATIONS:} \\
\texttt{1. "a red bird chase a stone"  =>  ketoge lupo pepa letu} \\
\texttt{2. "a bird chase a red stone"  =>  ketoge lupo letu pepa} \\
\texttt{3. "a small bird chase a stone"  =>  ketoge lupo pepa behito} \\
\texttt{4. "a bird chase a small stone"  =>  ketoge lupo behito pepa} \\
\texttt{STATUS: investigation round 2 of 5 done. Probes used: 8 of 20. 3 investigation round(s) left before the exam (10 sentences).} \\
\texttt{Send your next probes now (a PROBES list; brief THOUGHTS optional).} \\
            }
        }
    }
     \\ \\

    \theutterance \stepcounter{utterance}  
    & \multicolumn{4}{p{0.6\linewidth}}{
        \cellcolor[rgb]{0.561,0.737,0.561}{
            \makecell[{{p{\linewidth}}}]{
                \texttt{\tiny{[P1$\rangle$GM]}}
                \texttt{THOUGHTS: Excellent progress. I have confirmed the word order is SVO. I have identified the nouns: 'lupo' = bird, 'pepa' = stone. I have confirmed that adjectives come *after* the noun they modify (e.g., 'red bird' {-}> 'lupo letu', 'red stone' {-}> 'pepa letu'). I also see that the verb 'chase' is 'ketoge' and 'see' is 'pabi', 'hold' is 'patalo'. Now I need to test: 1) Do verbs change form based on the subject or object? 2) How does number work (pluralization)? 3) How does definiteness ('the' vs 'a') work? 4) Do adjectives agree in number? I will test pluralization with 'two birds' and 'two stones', and test definiteness with 'the bird'. I will also check if the verb changes when the subject is plural.} \\
\\ 
\texttt{PROBES:} \\
\texttt{1. two birds chase a stone} \\
\texttt{2. two stones chase a bird} \\
\texttt{3. the bird chase a stone} \\
\texttt{4. a bird chase the stone} \\
            }
        }
    }
    & & \\ \\

    \theutterance \stepcounter{utterance}  
    & & & \multicolumn{2}{p{0.3\linewidth}}{
        \cellcolor[rgb]{0.9,0.9,0.9}{
            \makecell[{{p{\linewidth}}}]{
                \texttt{\tiny{[GM$|$GM]}}
                {'kind': 'investigation', 'round': 2, 'probes_issued': 4, 'malformed': 0, 'probes': ['two birds chase a stone', 'two stones chase a bird', 'the bird chase a stone', 'a bird chase the stone']}
            }
        }
    }
    & & \\ \\

    \theutterance \stepcounter{utterance}  
    & & & \multicolumn{4}{p{0.6\linewidth}}{
        \cellcolor[rgb]{0.9,0.9,0.9}{
            \makecell[{{p{\linewidth}}}]{
                \texttt{\tiny{[P1$\langle$GM]}}
                \texttt{TRANSLATIONS:} \\
\texttt{1. "two birds chase a stone"  =>  ketoge lupo te pepahu} \\
\texttt{2. "two stones chase a bird"  =>  ketoge pepa te lupohu} \\
\texttt{3. "the bird chase a stone"  =>  ketoge lupo pepa la} \\
\texttt{4. "a bird chase the stone"  =>  ketoge lupo la pepa} \\
\texttt{STATUS: investigation round 3 of 5 done. Probes used: 12 of 20. 2 investigation round(s) left before the exam (10 sentences).} \\
\texttt{Send your next probes now (a PROBES list; brief THOUGHTS optional).} \\
            }
        }
    }
     \\ \\

    \theutterance \stepcounter{utterance}  
    & \multicolumn{4}{p{0.6\linewidth}}{
        \cellcolor[rgb]{0.561,0.737,0.561}{
            \makecell[{{p{\linewidth}}}]{
                \texttt{\tiny{[P1$\rangle$GM]}}
                \texttt{THOUGHTS: Great findings so far.} \\
\texttt{{-} Word order: SVO.} \\
\texttt{{-} Nouns: lupo (bird), pepa (stone).} \\
\texttt{{-} Adjectives: post{-}nominal (e.g., lupo letu).} \\
\texttt{{-} Number: Plural is marked by adding 'te' after the noun (lupo te) and a suffix 'hu' to the noun itself (pepahu, lupohu). So plural = Noun + te + NounSuffix.} \\
\texttt{{-} Definiteness: 'the' adds 'la' after the noun (lupo la).} \\
\texttt{{-} Verbs: So far, 'chase' is 'ketoge'. I need to check if verbs change based on the subject/object number or definiteness. I also need to verify the plural suffix for 'bird' (is it 'hu' like stone? The translation "two birds" was "lupo te pepahu" {-} wait, "pepahu" is the stone suffix. "lupo te" is the subject. The object is "a stone". Let's re{-}examine "two birds chase a stone": "ketoge lupo te pepahu". This looks like "chase {[}two birds{]} {[}stone{-}hu{]}". Wait, "pepahu" usually means plural stone. But the English is "a stone". Let me re{-}read carefully.} \\
\texttt{Ah, "two birds chase a stone" => "ketoge lupo te pepahu".} \\
\texttt{"two stones chase a bird" => "ketoge pepa te lupohu".} \\
\texttt{In the first case, "lupo te" is clearly "two birds". "pepahu" is "stone" + "hu". But the English object is singular "a stone". This implies "hu" might be a verb ending? Or maybe I misidentified the structure.} \\
\texttt{Let's look at "the bird chase a stone" => "ketoge lupo pepa la". Here "lupo" is singular bird, "pepa" is singular stone, "la" is "the".} \\
\texttt{In "two birds chase a stone", we have "lupo te" (two birds). The object is "pepahu". If "pepa" is stone, "pepahu" is stone+hu. But the object is singular. Maybe "hu" is a verb marker? No, the verb is at the start.} \\
\texttt{Alternative hypothesis: The word order in the plural sentence is different? No, "ketoge" is first.} \\
\texttt{Maybe "pepahu" is actually "a stone" in a specific context? No, that's unlikely.} \\
\texttt{Let's reconsider the plural marker.} \\
\texttt{"two birds" {-}> "lupo te".} \\
\texttt{"two stones" {-}> "pepa te".} \\
\texttt{In "two birds chase a stone", the object is "pepahu". If "pepa" is stone, then "hu" is attached. But the English says "a stone".} \\
\texttt{Wait, could "hu" be the verb for "chase"? No, "ketoge" is used.} \\
\texttt{Could "hu" be a marker for the *object* being singular? No, "pepa" is singular in "a bird chase a stone".} \\
\texttt{Let's look at "two stones chase a bird": "ketoge pepa te lupohu". Subject: "pepa te" (two stones). Object: "lupohu". "lupo" is bird. "hu" is attached. English object is "a bird".} \\
\texttt{Hypothesis: The suffix 'hu' marks the *object* of the sentence when the subject is plural? That seems weird.} \\
\texttt{Alternative: Maybe I misparsed "lupo te pepahu". Could it be "lupo" (bird) "te" (two) "pepa" (stone) "hu" (something)? No, "lupo te" is clearly the plural subject based on the second sentence.} \\
\texttt{Let's look at the singular sentences again.} \\
\texttt{"a bird chase a stone" {-}> "ketoge lupo pepa".} \\
\texttt{"the bird chase a stone" {-}> "ketoge lupo pepa la".} \\
\texttt{"a bird chase the stone" {-}> "ketoge lupo la pepa".} \\
\texttt{So "la" is definitely the definite article marker placed after the noun.} \\
\texttt{Now back to plurals.} \\
\texttt{"two birds chase a stone" {-}> "ketoge lupo te pepahu".} \\
\texttt{"two stones chase a bird" {-}> "ketoge pepa te lupohu".} \\
\texttt{In both cases, the *object* gets the suffix 'hu'. The subject gets 'te'.} \\
\texttt{Is it possible that 'hu' is a marker for "indefinite singular" in this specific construction? No, "pepa" is already indefinite singular.} \\
\texttt{Maybe 'hu' is a verb agreement marker for "plural subject"?} \\
\texttt{If Subject is plural, Verb + Object gets 'hu'?} \\
\texttt{Let's test this with a new sentence where the subject is singular but the object is plural.} \\
\texttt{"a bird chase two stones".} \\
\texttt{Expected: Subject "a bird" (lupo), Verb "chase" (ketoge), Object "two stones" (pepa te).} \\
\texttt{If the rule is "Plural Subject {-}> Object gets hu", then "a bird" (singular) should NOT have 'hu'.} \\
\texttt{If the rule is "Plural Object {-}> Object gets te + hu", then "two stones" should be "pepa te hu"?} \\
\texttt{Let's try to construct "a bird chase two stones".} \\
\texttt{If my previous analysis of "two birds chase a stone" is correct, the object "a stone" became "pepahu". This is very strange unless "hu" means "one" or "a" in a specific grammatical slot, or maybe "pepahu" is the word for "a stone" and "pepa" is just "stone"? No, "a stone" was "pepa".} \\
\texttt{Wait, let's re{-}read the first plural sentence: "two birds chase a stone" => "ketoge lupo te pepahu".} \\
\texttt{Maybe the word order is Subject{-}Verb{-}Object{-}Number?} \\
\texttt{"two birds" = "lupo te".} \\
\texttt{"a stone" = "pepa hu"?} \\
\texttt{If "hu" marks indefinite singular object?} \\
\texttt{Let's check "two stones chase a bird" => "ketoge pepa te lupohu".} \\
\texttt{"two stones" = "pepa te".} \\
\texttt{"a bird" = "lupo hu".} \\
\texttt{This pattern holds: When the object is indefinite singular, it gets 'hu'.} \\
\texttt{But in "a bird chase a stone" => "ketoge lupo pepa", there is no 'hu'.} \\
\texttt{Why the difference?} \\
\texttt{Maybe the verb changes? No, "ketoge" is constant.} \\
\texttt{Maybe the position of 'te' changes?} \\
\texttt{Let's reconsider the "two birds" phrase. "lupo te".} \\
\texttt{Is it possible that "lupo te" is actually "bird two" and the order is Noun{-}Number?} \\
\texttt{Then "two birds chase a stone" {-}> "lupo te ketoge pepa hu"? No, the translation is "ketoge lupo te pepahu".} \\
\texttt{Okay, let's look at the structure: {[}Verb{]} {[}Subject{]} {[}Object{]}.} \\
\texttt{Sentence 1: ketoge (V) lupo te (S) pepahu (O).} \\
\texttt{Sentence 2: ketoge (V) pepa te (S) lupohu (O).} \\
\texttt{In both, the Object has 'hu'. The Subject has 'te'.} \\
\texttt{In singular sentences:} \\
\texttt{"a bird chase a stone" {-}> ketoge lupo pepa. (No 'te', no 'hu').} \\
\texttt{"the bird chase a stone" {-}> ketoge lupo pepa la. (No 'te', 'la' added).} \\
\texttt{"a bird chase the stone" {-}> ketoge lupo la pepa. ('la' added).} \\
\texttt{Hypothesis:} \\
\texttt{{-} Plural Subject: Add 'te' after the noun.} \\
\texttt{{-} Indefinite Singular Object: Add 'hu' after the noun?} \\
\texttt{But why didn't "a bird chase a stone" have 'hu'?} \\
\texttt{Maybe 'hu' is only added when the *subject* is plural?} \\
\texttt{Let's test: "a bird chase two stones".} \\
\texttt{Subject: "a bird" (lupo). Object: "two stones" (pepa te).} \\
\texttt{If the rule is "Plural Subject {-}> Object gets hu", then this sentence should NOT have 'hu'.} \\
\texttt{If the rule is "Plural Object {-}> Object gets te + hu", then "two stones" might be "pepa te hu".} \\
\texttt{Let's try to generate "a bird chase two stones".} \\
\texttt{If the pattern "Plural Subject {-}> Object gets hu" is true, then "a bird chase two stones" should be "ketoge lupo pepa te".} \\
\texttt{If the pattern "Plural Object {-}> Object gets te + hu" is true, then "two stones" is "pepa te hu", so "ketoge lupo pepa te hu".} \\
\texttt{Let's also test "two birds chase two stones".} \\
\texttt{Subject: "two birds" (lupo te). Object: "two stones" (pepa te).} \\
\texttt{If "Plural Subject {-}> Object gets hu", then "ketoge lupo te pepa te hu".} \\
\texttt{If "Plural Object {-}> Object gets te + hu", then "ketoge lupo te pepa te hu".} \\
\texttt{It seems "te" marks plural count, and "hu" marks something else.} \\
\texttt{Wait, look at "two birds chase a stone" {-}> "lupo te pepahu".} \\
\texttt{"two stones chase a bird" {-}> "pepa te lupohu".} \\
\texttt{In both cases, the *object* is indefinite singular ("a stone", "a bird") and it gets 'hu'.} \\
\texttt{In "a bird chase a stone", the object is indefinite singular ("a stone") and it does NOT get 'hu'.} \\
\texttt{The only difference is the Subject: "two birds" (plural) vs "a bird" (singular).} \\
\texttt{Conclusion: When the Subject is plural, the Indefinite Singular Object gets the suffix 'hu'.} \\
\texttt{What about the Definite Singular Object?} \\
\texttt{"a bird chase the stone" {-}> "ketoge lupo la pepa". No 'hu'.} \\
\texttt{"two birds chase the stone"?} \\
\texttt{Subject: "two birds" (lupo te). Object: "the stone" (pepa la).} \\
\texttt{Prediction: Should it be "ketoge lupo te pepa la hu"? Or does 'la' block 'hu'? Or does 'hu' only apply to indefinite?} \\
\texttt{Let's test "two birds chase the stone".} \\
\texttt{Also, I need to check if the verb changes. "chase" is "ketoge". Does it change for plural subject?} \\
\texttt{Let's try "two birds chase a stone" again. It was "ketoge".} \\
\texttt{Let's try "a bird chase two stones".} \\
\texttt{If the rule is "Plural Subject {-}> Object gets hu", then "a bird" (singular subject) {-}> no 'hu'.} \\
\texttt{So "a bird chase two stones" {-}> "ketoge lupo pepa te".} \\
\texttt{If the rule is "Plural Object {-}> Object gets te + hu", then "two stones" {-}> "pepa te hu".} \\
\texttt{So "a bird chase two stones" {-}> "ketoge lupo pepa te hu".} \\
\texttt{This is a crucial distinction. I need to know if 'hu' is triggered by the *subject* being plural or the *object* being plural.} \\
\texttt{Given "two birds chase a stone" {-}> "pepahu" (object indefinite singular + hu), it strongly suggests 'hu' is triggered by the *subject* being plural.} \\
\texttt{But wait, if 'hu' is triggered by plural subject, why isn't it in "a bird chase a stone"? Correct.} \\
\texttt{Why isn't it in "a bird chase the stone"? Correct.} \\
\texttt{So 'hu' = "Plural Subject Marker on Indefinite Singular Object"?} \\
\texttt{Or maybe 'hu' is just a marker for "Indefinite Singular Object" but it's omitted in some contexts? No, that's unlikely.} \\
\texttt{Let's assume:} \\
\texttt{1. Plural Subject: Noun + te.} \\
\texttt{2. Indefinite Singular Object: Noun + hu (ONLY IF Subject is plural).} \\
\texttt{3. Plural Object: Noun + te + hu? Or just Noun + te?} \\
\texttt{Let's test "a bird chase two stones".} \\
\texttt{If Plural Object = Noun + te, then "ketoge lupo pepa te".} \\
\texttt{If Plural Object = Noun + te + hu, then "ketoge lupo pepa te hu".} \\
\texttt{Let's also test "two birds chase two stones".} \\
\texttt{If Plural Subject + Plural Object: "ketoge lupo te pepa te hu"?} \\
\texttt{Let's try to differentiate.} \\
\texttt{I will ask:} \\
\texttt{1. "a bird chase two stones" (Singular Subject, Plural Object)} \\
\texttt{2. "two birds chase two stones" (Plural Subject, Plural Object)} \\
\texttt{3. "two birds chase the stone" (Plural Subject, Definite Object) {-} to see if 'hu' applies to definite objects too.} \\
\\ 
\texttt{PROBES:} \\
\texttt{1. a bird chase two stones} \\
\texttt{2. two birds chase two stones} \\
\texttt{3. two birds chase the stone} \\
            }
        }
    }
    & & \\ \\

    \theutterance \stepcounter{utterance}  
    & & & \multicolumn{2}{p{0.3\linewidth}}{
        \cellcolor[rgb]{0.9,0.9,0.9}{
            \makecell[{{p{\linewidth}}}]{
                \texttt{\tiny{[GM$|$GM]}}
                {'kind': 'investigation', 'round': 3, 'probes_issued': 3, 'malformed': 0, 'probes': ['a bird chase two stones', 'two birds chase two stones', 'two birds chase the stone']}
            }
        }
    }
    & & \\ \\

    \theutterance \stepcounter{utterance}  
    & & & \multicolumn{4}{p{0.6\linewidth}}{
        \cellcolor[rgb]{0.9,0.9,0.9}{
            \makecell[{{p{\linewidth}}}]{
                \texttt{\tiny{[P1$\langle$GM]}}
                \texttt{TRANSLATIONS:} \\
\texttt{1. "a bird chase two stones"  =>  ketoge te lupohu pepa} \\
\texttt{2. "two birds chase two stones"  =>  ketoge te lupohu te pepahu} \\
\texttt{3. "two birds chase the stone"  =>  ketoge lupo la te pepahu} \\
\texttt{STATUS: investigation round 4 of 5 done. Probes used: 15 of 20. 1 investigation round(s) left before the exam (10 sentences).} \\
\texttt{Send your next probes now (a PROBES list; brief THOUGHTS optional).} \\
            }
        }
    }
     \\ \\

    \theutterance \stepcounter{utterance}  
    & \multicolumn{4}{p{0.6\linewidth}}{
        \cellcolor[rgb]{0.561,0.737,0.561}{
            \makecell[{{p{\linewidth}}}]{
                \texttt{\tiny{[P1$\rangle$GM]}}
                \texttt{THOUGHTS: The previous results were confusing because I misparsed the word order. Let's re{-}analyze carefully:} \\
\texttt{1. "a bird chase two stones" => `ketoge te lupohu pepa`} \\
\texttt{   {-} Structure: Verb + {[}te + lupohu{]} + {[}pepa{]}} \\
\texttt{   {-} "te lupohu" corresponds to "two stones". So "te" is the number word, and "lupohu" is the plural noun "stones".} \\
\texttt{   {-} "pepa" is the singular noun "stone".} \\
\texttt{   {-} Order: Verb {-} Object {-} Subject? Or Verb {-} (Number+Noun) {-} (Noun)?} \\
\texttt{   {-} Wait, "two stones" is the object. "a stone" is the subject? No, "a bird" is the subject.} \\
\texttt{   {-} If "a bird" is the subject, where is it? It's missing or I misidentified "lupo".} \\
\texttt{   {-} Let's look at "two birds chase a stone" => `ketoge lupo te pepahu`.} \\
\texttt{     {-} Here "lupo te" looks like "two birds". "pepahu" looks like "a stone".} \\
\texttt{     {-} But in the new sentence "a bird chase two stones" => `ketoge te lupohu pepa`.} \\
\texttt{     {-} If "te lupohu" is "two stones", then "pepa" must be "a bird".} \\
\texttt{     {-} This contradicts my earlier finding that "lupo" = bird and "pepa" = stone.} \\
\texttt{     {-} Let's re{-}evaluate the nouns.} \\
\texttt{     {-} Sentence 1: "a bird chase a stone" => `ketoge lupo pepa`.} \\
\texttt{       {-} If "lupo" = bird, "pepa" = stone.} \\
\texttt{     {-} Sentence 5: "a bird chase two stones" => `ketoge te lupohu pepa`.} \\
\texttt{       {-} If "pepa" = stone (singular), then "te lupohu" must be "two stones"? No, "two stones" should be plural.} \\
\texttt{       {-} If "te lupohu" = "two stones", then "lupohu" = stone (plural).} \\
\texttt{       {-} Then "pepa" = "a bird".} \\
\texttt{       {-} This implies "pepa" = bird and "lupo" = stone?} \\
\texttt{       {-} Let's check Sentence 2: "two stones chase a bird" => `ketoge pepa te lupohu`.} \\
\texttt{         {-} If "pepa" = bird, "lupohu" = stone (plural).} \\
\texttt{         {-} "pepa te" = "two birds"?} \\
\texttt{         {-} "lupohu" = "a stone"?} \\
\texttt{         {-} This fits if: "pepa" = bird, "lupo" = stone.} \\
\texttt{         {-} Let's re{-}verify with Sentence 1: "a bird chase a stone" => `ketoge lupo pepa`.} \\
\texttt{           {-} If "lupo" = stone, "pepa" = bird.} \\
\texttt{           {-} "ketoge stone bird"? That would mean "stone chase bird". But English is "bird chase stone".} \\
\texttt{           {-} So the order is Verb {-} Object {-} Subject?} \\
\texttt{           {-} "ketoge" (chase) {-} "lupo" (stone) {-} "pepa" (bird). {-}> "Stone chase bird". Incorrect meaning.} \\
\texttt{           {-} Maybe the order is Verb {-} Subject {-} Object, but I swapped the nouns.} \\
\texttt{           {-} If "lupo" = bird, "pepa" = stone.} \\
\texttt{           {-} "a bird chase a stone" => `ketoge lupo pepa`. (V S O). Correct.} \\
\texttt{           {-} "a bird chase two stones" => `ketoge te lupohu pepa`.} \\
\texttt{             {-} If V S O, then "te lupohu" is Subject ("two birds"?) and "pepa" is Object ("a stone"?).} \\
\texttt{             {-} But English is "a bird chase two stones". Subject="a bird", Object="two stones".} \\
\texttt{             {-} So "te lupohu" must be "two stones" and "pepa" must be "a bird".} \\
\texttt{             {-} This implies "pepa" = bird, "lupo" = stone.} \\
\texttt{             {-} But then "a bird chase a stone" => `ketoge lupo pepa` would be "chase stone bird" (V O S?).} \\
\texttt{             {-} If V O S: "chase stone bird" = "Bird chases stone".} \\
\texttt{             {-} Let's check "two stones chase a bird" => `ketoge pepa te lupohu`.} \\
\texttt{               {-} If V O S: "chase {[}two birds{]} {[}a stone{]}"? No.} \\
\texttt{               {-} If V O S: "chase {[}pepa te{]} {[}lupohu{]}".} \\
\texttt{               {-} If "pepa" = bird, "lupo" = stone.} \\
\texttt{               {-} "pepa te" = "two birds". "lupohu" = "a stone".} \\
\texttt{               {-} Translation: "chase two birds a stone". Meaning: "Two birds chase a stone"? No, English is "Two stones chase a bird".} \\
\texttt{               {-} This is getting contradictory. Let's look at the "te" position.} \\
\texttt{               {-} In "two birds chase a stone" => `ketoge lupo te pepahu`.} \\
\texttt{                 {-} "lupo te" = "two birds"? "pepahu" = "a stone"?} \\
\texttt{               {-} In "a bird chase two stones" => `ketoge te lupohu pepa`.} \\
\texttt{                 {-} "te lupohu" = "two stones"? "pepa" = "a bird"?} \\
\texttt{                 {-} If "lupo" = bird, "pepa" = stone.} \\
\texttt{                 {-} Then "lupo te" = "two birds". "pepahu" = "a stone".} \\
\texttt{                 {-} And "te lupohu" = "two stones"? "lupohu" = "stone"?} \\
\texttt{                 {-} This implies "lupo" and "lupohu" are related to stone/bird.} \\
\texttt{                 {-} Maybe "lupo" = bird, "pepa" = stone.} \\
\texttt{                 {-} "two birds" = "lupo te".} \\
\texttt{                 {-} "two stones" = "pepa te".} \\
\texttt{                 {-} "a stone" = "pepa".} \\
\texttt{                 {-} "a bird" = "lupo".} \\
\texttt{                 {-} Let's re{-}parse "a bird chase two stones" => `ketoge te lupohu pepa`.} \\
\texttt{                   {-} If V S O: "chase" "te lupohu" "pepa".} \\
\texttt{                   {-} "te lupohu" = "two stones"? "pepa" = "a bird"?} \\
\texttt{                   {-} This requires "pepa" = bird.} \\
\texttt{                   {-} But "a bird chase a stone" => `ketoge lupo pepa`.} \\
\texttt{                   {-} If "pepa" = bird, then "lupo" = stone.} \\
\texttt{                   {-} "chase" "stone" "bird". (V O S).} \\
\texttt{                   {-} Does "stone chase bird" mean "bird chases stone"? No.} \\
\texttt{                   {-} Maybe the language is V{-}O{-}S and the meaning is reversed? No, "chase" is transitive.} \\
\texttt{                   {-} Let's reconsider the noun mapping.} \\
\texttt{                   {-} Maybe "lupo" = stone, "pepa" = bird.} \\
\texttt{                   {-} "a bird chase a stone" => `ketoge lupo pepa`.} \\
\texttt{                     {-} V O S: "chase stone bird". (Bird chases stone).} \\
\texttt{                   {-} "a bird chase two stones" => `ketoge te lupohu pepa`.} \\
\texttt{                     {-} V O S: "chase {[}two stones{]} bird".} \\
\texttt{                     {-} "te lupohu" = "two stones". "pepa" = "bird".} \\
\texttt{                     {-} This works if "lupohu" = stone (plural) and "pepa" = bird.} \\
\texttt{                   {-} "two stones chase a bird" => `ketoge pepa te lupohu`.} \\
\texttt{                     {-} V O S: "chase {[}two birds{]} stone"? No.} \\
\texttt{                     {-} V O S: "chase {[}pepa te{]} {[}lupohu{]}".} \\
\texttt{                     {-} If "pepa" = bird, "lupohu" = stone.} \\
\texttt{                     {-} "chase {[}two birds{]} {[}stone{]}". (Two birds chase a stone).} \\
\texttt{                     {-} But English is "Two stones chase a bird".} \\
\texttt{                     {-} So V O S doesn't fit all.} \\
\\ 
\texttt{Let's try V{-}S{-}O again with a different noun mapping.} \\
\texttt{Maybe "lupo" = bird, "pepa" = stone.} \\
\texttt{"a bird chase a stone" => `ketoge lupo pepa`. (V S O).} \\
\texttt{"a bird chase two stones" => `ketoge te lupohu pepa`.} \\
\texttt{{-} If V S O: "chase" "te lupohu" "pepa".} \\
\texttt{{-} "te lupohu" = "two stones"? "pepa" = "a bird"?} \\
\texttt{{-} This implies "pepa" = bird. Contradiction.} \\
\\ 
\texttt{Is it possible the order is V{-}O{-}S and the nouns are:} \\
\texttt{"lupo" = stone, "pepa" = bird.} \\
\texttt{"a bird chase a stone" => `ketoge lupo pepa`. (V O S).} \\
\texttt{"a bird chase two stones" => `ketoge te lupohu pepa`.} \\
\texttt{{-} V O S: "chase" "te lupohu" "pepa".} \\
\texttt{{-} "te lupohu" = "two stones". "pepa" = "bird".} \\
\texttt{{-} "lupohu" = stone (plural). "pepa" = bird.} \\
\texttt{{-} "two stones chase a bird" => `ketoge pepa te lupohu`.} \\
\texttt{{-} V O S: "chase" "pepa te" "lupohu".} \\
\texttt{{-} "pepa te" = "two birds". "lupohu" = "stone".} \\
\texttt{{-} Translation: "chase two birds a stone". (Two birds chase a stone).} \\
\texttt{{-} But English is "Two stones chase a bird".} \\
\texttt{{-} So this mapping fails.} \\
\\ 
\texttt{Let's look at the "te" position again.} \\
\texttt{In "two birds chase a stone" => `ketoge lupo te pepahu`.} \\
\texttt{{-} "lupo te" = "two birds".} \\
\texttt{{-} "pepahu" = "a stone".} \\
\texttt{In "a bird chase two stones" => `ketoge te lupohu pepa`.} \\
\texttt{{-} "te lupohu" = "two stones".} \\
\texttt{{-} "pepa" = "a bird".} \\
\texttt{This implies "lupo" = bird, "pepa" = stone.} \\
\texttt{And "pepahu" = stone (plural)? "lupohu" = bird (plural)?} \\
\texttt{If "pepahu" = stone (plural), then "a stone" (singular) is "pepa".} \\
\texttt{If "lupohu" = bird (plural), then "a bird" (singular) is "lupo".} \\
\texttt{Let's check "two stones chase a bird" => `ketoge pepa te lupohu`.} \\
\texttt{{-} If "pepa" = stone, "lupohu" = bird (plural).} \\
\texttt{{-} "pepa te" = "two stones".} \\
\texttt{{-} "lupohu" = "a bird"? No, "lupohu" is plural form.} \\
\texttt{{-} So "two stones chase a bird" => "chase" "two stones" "a bird".} \\
\texttt{{-} Structure: V O S?} \\
\texttt{{-} "ketoge" (V) "pepa te" (O: two stones) "lupohu" (S: a bird??).} \\
\texttt{{-} But "lupohu" is plural. "a bird" is singular.} \\
\texttt{{-} Unless "lupohu" means "a bird" in this context? No, "hu" seems to be a plural marker.} \\
\texttt{{-} Wait, in "two birds chase a stone" => `ketoge lupo te pepahu`.} \\
\texttt{{-} "lupo te" = "two birds". "pepahu" = "a stone".} \\
\texttt{{-} Here "pepahu" is singular object.} \\
\texttt{{-} In "a bird chase two stones" => `ketoge te lupohu pepa`.} \\
\texttt{{-} "te lupohu" = "two stones". "pepa" = "a bird".} \\
\texttt{{-} Here "pepa" is singular object.} \\
\texttt{{-} In "two stones chase a bird" => `ketoge pepa te lupohu`.} \\
\texttt{{-} "pepa te" = "two stones". "lupohu" = "a bird".} \\
\texttt{{-} Here "lupohu" is singular object? But "hu" is plural marker?} \\
\texttt{{-} This is very confusing. Let's reconsider the "hu" suffix.} \\
\texttt{{-} Maybe "hu" is not plural. Maybe "hu" is a marker for the *object* when the subject is plural?} \\
\texttt{{-} "two birds chase a stone" (S=Plural) => Object "a stone" gets "hu" {-}> "pepahu".} \\
\texttt{{-} "two stones chase a bird" (S=Plural) => Object "a bird" gets "hu" {-}> "lupohu".} \\
\texttt{{-} "a bird chase two stones" (S=Singular) => Object "two stones" does NOT get "hu".} \\
\texttt{{-} "a bird chase a stone" (S=Singular) => Object "a stone" does NOT get "hu".} \\
\texttt{{-} This fits perfectly!} \\
\texttt{{-} So:} \\
\texttt{  {-} "hu" = Object marker for Plural Subject.} \\
\texttt{  {-} "te" = Number word "two" (or plural marker?).} \\
\texttt{  {-} In "two birds chase a stone": "lupo te" = "two birds". "pepahu" = "a stone" + hu.} \\
\texttt{  {-} In "a bird chase two stones": "te lupohu" = "two stones". "pepa" = "a bird".} \\
\texttt{    {-} Wait, "te lupohu". If "te" is "two", then "lupohu" is "stones".} \\
\texttt{    {-} But "lupo" is "bird". So "lupohu" = "bird" + "hu"?} \\
\texttt{    {-} If "lupohu" = "stones", then "lupo" must be "stone".} \\
\texttt{    {-} Let's swap nouns: "lupo" = stone, "pepa" = bird.} \\
\texttt{    {-} "a bird chase a stone" => `ketoge lupo pepa`.} \\
\texttt{      {-} V S O: "chase" "stone" "bird". (Bird chases stone).} \\
\texttt{      {-} "lupo" = stone. "pepa" = bird.} \\
\texttt{    {-} "two birds chase a stone" => `ketoge lupo te pepahu`.} \\
\texttt{      {-} V S O: "chase" "stone two" "bird{-}hu"?} \\
\texttt{      {-} "lupo te" = "stone two"? No.} \\
\texttt{      {-} "lupo" = stone. "te" = two. "pepahu" = bird{-}hu?} \\
\texttt{      {-} This doesn't fit "two birds".} \\
\\ 
\texttt{Let's go back to: "lupo" = bird, "pepa" = stone.} \\
\texttt{"a bird chase a stone" => `ketoge lupo pepa`. (V S O).} \\
\texttt{"two birds chase a stone" => `ketoge lupo te pepahu`.} \\
\texttt{{-} "lupo te" = "two birds".} \\
\texttt{{-} "pepahu" = "a stone" + hu.} \\
\texttt{"a bird chase two stones" => `ketoge te lupohu pepa`.} \\
\texttt{{-} "te lupohu" = "two stones".} \\
\texttt{{-} "pepa" = "a bird".} \\
\texttt{{-} This implies "lupohu" = "stones".} \\
\texttt{{-} If "lupo" = bird, then "lupohu" = bird + hu?} \\
\texttt{{-} But "te lupohu" = "two stones".} \\
\texttt{{-} So "lupohu" must mean "stones".} \\
\texttt{{-} This implies "lupo" = stone?} \\
\texttt{{-} Contradiction.} \\
\\ 
\texttt{Alternative: The word order is V{-}O{-}S.} \\
\texttt{"a bird chase a stone" => `ketoge lupo pepa`.} \\
\texttt{{-} V O S: "chase" "bird" "stone". (Bird chases stone).} \\
\texttt{{-} "lupo" = bird. "pepa" = stone.} \\
\texttt{"two birds chase a stone" => `ketoge lupo te pepahu`.} \\
\texttt{{-} V O S: "chase" "bird two" "stone{-}hu"?} \\
\texttt{{-} "lupo te" = "two birds".} \\
\texttt{{-} "pepahu" = "a stone" + hu.} \\
\texttt{"a bird chase two stones" => `ketoge te lupohu pepa`.} \\
\texttt{{-} V O S: "chase" "two stone{-}hu" "bird"?} \\
\texttt{{-} "te lupohu" = "two stones".} \\
\texttt{{-} "pepa" = "a bird".} \\
\texttt{{-} This implies "lupohu" = "stones".} \\
\texttt{{-} If "lupo" = bird, then "lupohu" = bird + hu?} \\
\texttt{{-} But "te lupohu" = "two stones".} \\
\texttt{{-} So "lupohu" must be "stones".} \\
\texttt{{-} This implies "lupo" = stone.} \\
\texttt{{-} But "lupo" was "bird" in the first sentence.} \\
\\ 
\texttt{Is it possible that "lupo" and "pepa" are not the nouns I think they are?} \\
\texttt{Maybe "lupo" = stone, "pepa" = bird.} \\
\texttt{"a bird chase a stone" => `ketoge lupo pepa`.} \\
\texttt{{-} V O S: "chase" "stone" "bird". (Bird chases stone).} \\
\texttt{{-} "lupo" = stone. "pepa" = bird.} \\
\texttt{"two birds chase a stone" => `ketoge lupo te pepahu`.} \\
\texttt{{-} V O S: "chase" "stone two" "bird{-}hu"?} \\
\texttt{{-} "lupo te" = "stone two"? No.} \\
\texttt{{-} "lupo" = stone. "te" = two. "pepahu" = bird{-}hu?} \\
\texttt{{-} "two birds" = "lupo te"? No, "lupo" is stone.} \\
\texttt{{-} Maybe "te" is the plural marker for the noun before it?} \\
\texttt{{-} "lupo te" = "two stones"?} \\
\texttt{{-} Then "pepahu" = "a bird" + hu.} \\
\texttt{{-} "two stones chase a bird" => `ketoge pepa te lupohu`.} \\
\texttt{{-} V O S: "chase" "bird two" "stone{-}hu"?} \\
\texttt{{-} "pepa te" = "two birds".} \\
\texttt{{-} "lupohu" = "a stone" + hu.} \\
\texttt{{-} This fits if:} \\
\texttt{  {-} "lupo" = stone.} \\
\texttt{  {-} "pepa" = bird.} \\
\texttt{  {-} "te" = "two" (or plural marker).} \\
\texttt{  {-} "hu" = marker for object when subject is plural?} \\
\texttt{  {-} Let's check "a bird chase two stones" => `ketoge te lupohu pepa`.} \\
\texttt{    {-} V O S: "chase" "two stone{-}hu" "bird"?} \\
\texttt{    {-} "te lupohu" = "two stones".} \\
\texttt{    {-} "pepa" = "a bird".} \\
\texttt{    {-} "lupohu" = "stones".} \\
\texttt{    {-} If "lupo" = stone, then "lupohu" = stone + hu?} \\
\texttt{    {-} But "te lupohu" = "two stones".} \\
\texttt{    {-} So "lupohu" = "stones".} \\
\texttt{    {-} "pepa" = "bird".} \\
\texttt{    {-} "two birds" = "pepa te".} \\
\texttt{    {-} "two stones" = "lupo te"?} \\
\texttt{    {-} Let's check "two birds chase a stone" => `ketoge lupo te pepahu`.} \\
\texttt{      {-} V O S: "chase" "stone two" "bird{-}hu"?} \\
\texttt{      {-} "lupo te" = "two stones".} \\
\texttt{      {-} "pepahu" = "a bird" + hu.} \\
\texttt{      {-} But English is "two birds chase a stone".} \\
\texttt{      {-} So "lupo te" must be "two birds".} \\
\texttt{      {-} This implies "lupo" = bird.} \\
\texttt{      {-} But "lupo" = stone in the previous hypothesis.} \\
\\ 
\texttt{Okay, let's try one more time with V{-}S{-}O and "lupo"=bird, "pepa"=stone.} \\
\texttt{1. "a bird chase a stone" => `ketoge lupo pepa`. (V S O).} \\
\texttt{2. "two birds chase a stone" => `ketoge lupo te pepahu`.} \\
\texttt{   {-} "lupo te" = "two birds".} \\
\texttt{   {-} "pepahu" = "a stone" + hu.} \\
\texttt{3. "a bird chase two stones" => `ketoge te lupohu pepa`.} \\
\texttt{   {-} "te lupohu" = "two stones".} \\
\texttt{   {-} "pepa" = "a bird".} \\
\texttt{   {-} This implies "lupohu" = "stones".} \\
\texttt{   {-} If "lupo" = bird, then "lupohu" = bird + hu?} \\
\texttt{   {-} But "te lupohu" = "two stones".} \\
\texttt{   {-} So "lupohu" must be "stones".} \\
\texttt{   {-} This implies "lupo" = stone.} \\
\texttt{   {-} Contradiction.} \\
\\ 
\texttt{Wait, what if "te" is not "two"? What if "te" is a plural marker?} \\
\texttt{"two birds" = "lupo te".} \\
\texttt{"two stones" = "pepa te".} \\
\texttt{"a bird" = "lupo".} \\
\texttt{"a stone" = "pepa".} \\
\texttt{"two birds chase a stone" => `ketoge lupo te pepahu`.} \\
\texttt{{-} "lupo te" = "two birds".} \\
\texttt{{-} "pepahu" = "a stone" + hu.} \\
\texttt{"a bird chase two stones" => `ketoge te lupohu pepa`.} \\
\texttt{{-} "te lupohu" = "two stones".} \\
\texttt{{-} "pepa" = "a bird".} \\
\texttt{{-} This implies "lupohu" = "stones".} \\
\texttt{{-} If "pepa" = stone, then "lupohu" = stone + hu?} \\
\texttt{{-} But "te lupohu" = "two stones".} \\
\texttt{{-} So "lupohu" = "stones".} \\
\texttt{{-} "pepa" = "bird".} \\
\texttt{{-} But "lupo" = "bird" in sentence 1.} \\
\texttt{{-} This is a hard contradiction unless "lupo" and "pepa" are not the nouns.} \\
\\ 
\texttt{Let's reconsider the words.} \\
\texttt{Maybe "lupo" = stone, "pepa" = bird.} \\
\texttt{1. "a bird chase a stone" => `ketoge lupo pepa`.} \\
\texttt{   {-} V S O: "chase" "stone" "bird". (Bird chases stone).} \\
\texttt{   {-} "lupo" = stone. "pepa" = bird.} \\
\texttt{2. "two birds chase a stone" => `ket} \\
            }
        }
    }
    & & \\ \\

    \theutterance \stepcounter{utterance}  
    & & & \multicolumn{2}{p{0.3\linewidth}}{
        \cellcolor[rgb]{0.9,0.9,0.9}{
            \makecell[{{p{\linewidth}}}]{
                \texttt{\tiny{[GM$|$GM]}}
                \texttt{could not find any probe sentences; put them under a 'PROBES:' line, one per line} \\
            }
        }
    }
    & & \\ \\

    \theutterance \stepcounter{utterance}  
    & & & \multicolumn{4}{p{0.6\linewidth}}{
        \cellcolor[rgb]{0.9,0.9,0.9}{
            \makecell[{{p{\linewidth}}}]{
                \texttt{\tiny{[P1$\langle$GM]}}
                \texttt{Your last response could not be processed: could not find any probe sentences; put them under a 'PROBES:' line, one per line.} \\
\texttt{List your probes (a brief THOUGHTS line is optional):} \\
\texttt{PROBES:} \\
\texttt{1. <an English sentence>} \\
\texttt{(up to 4 probe lines; at least one)} \\
            }
        }
    }
     \\ \\

    \theutterance \stepcounter{utterance}  
    & \multicolumn{4}{p{0.6\linewidth}}{
        \cellcolor[rgb]{0.561,0.737,0.561}{
            \makecell[{{p{\linewidth}}}]{
                \texttt{\tiny{[P1$\rangle$GM]}}
                \texttt{THOUGHTS: I need to resolve the noun confusion and the "hu" suffix rule. My previous analysis was tangled. Let's restart the noun mapping with fresh eyes based on the most consistent data point: "a bird chase a stone" = `ketoge lupo pepa`. Assuming SVO (Subject{-}Verb{-}Object) is standard, `lupo` is likely the subject (bird) and `pepa` is the object (stone).} \\
\texttt{If `lupo` = bird and `pepa` = stone:} \\
\texttt{{-} "two birds chase a stone" = `ketoge lupo te pepahu`. Here `lupo te` = "two birds". `pepahu` = "a stone" + `hu`.} \\
\texttt{{-} "a bird chase two stones" = `ketoge te lupohu pepa`. Here `te lupohu` = "two stones". `pepa` = "a bird".} \\
\texttt{This implies `lupohu` = "stones" and `pepa` = "bird". This contradicts `lupo`=bird.} \\
\texttt{Hypothesis: The word order is V{-}O{-}S (Verb{-}Object{-}Subject).} \\
\texttt{{-} "a bird chase a stone" = `ketoge lupo pepa`. V=`ketoge`, O=`lupo` (stone?), S=`pepa` (bird?).} \\
\texttt{{-} "two birds chase a stone" = `ketoge lupo te pepahu`. V=`ketoge`, O=`lupo te` (two stones?), S=`pepahu` (bird + hu?).} \\
\texttt{{-} "a bird chase two stones" = `ketoge te lupohu pepa`. V=`ketoge`, O=`te lupohu` (two birds?), S=`pepa` (stone?).} \\
\texttt{This fits perfectly if:} \\
\texttt{{-} `lupo` = stone.} \\
\texttt{{-} `pepa` = bird.} \\
\texttt{{-} `te` = "two" (or plural marker).} \\
\texttt{{-} `hu` = suffix added to the Subject when the Subject is plural? Or maybe `hu` marks the subject in a specific way?} \\
\texttt{Let's check "two stones chase a bird" = `ketoge pepa te lupohu`.} \\
\texttt{{-} V=`ketoge`.} \\
\texttt{{-} O=`pepa te` (two birds?).} \\
\texttt{{-} S=`lupohu` (stone + hu?).} \\
\texttt{If O=`pepa te` (two birds), then `pepa`=bird.} \\
\texttt{If S=`lupohu` (stone), then `lupo`=stone.} \\
\texttt{This confirms: `lupo` = stone, `pepa` = bird.} \\
\texttt{Order: V{-}O{-}S.} \\
\texttt{Now, what is `hu`?} \\
\texttt{{-} "two birds chase a stone" (S=Plural) {-}> S=`pepahu` (bird+hu).} \\
\texttt{{-} "a bird chase two stones" (S=Singular) {-}> S=`pepa` (bird). No `hu`.} \\
\texttt{{-} "two stones chase a bird" (S=Plural) {-}> S=`lupohu` (stone+hu).} \\
\texttt{{-} "a bird chase a stone" (S=Singular) {-}> S=`pepa` (bird). No `hu`.} \\
\texttt{Rule: If the Subject is plural, add `hu` to the Subject noun.} \\
\texttt{Also, where is the number word?} \\
\texttt{{-} "two birds" = `pepa te` (in V{-}O{-}S order: O=`pepa te`). Wait, in "two birds chase a stone", the object is "a stone" (`lupo`). The subject is "two birds" (`pepa te`).} \\
\texttt{{-} Translation: `ketoge lupo te pepahu`.} \\
\texttt{{-} V=`ketoge`. O=`lupo` (stone). S=`pepa te` (two birds) + `hu`? No, `pepahu` is one word.} \\
\texttt{{-} So `pepa` (bird) + `te` (two) + `hu` (plural subject marker) = `pepahu`? No, `te` is between `pepa` and `hu`?} \\
\texttt{{-} `pepa` + `te` + `hu` {-}> `pepahu`? No, `te` is clearly "two".} \\
\texttt{{-} Maybe `te` is the number, and `hu` is the plural marker.} \\
\texttt{{-} "two birds" = `pepa te`. With plural subject marker `hu` {-}> `pepa te hu`? But the output is `pepahu`.} \\
\texttt{{-} Maybe `te` and `hu` merge? Or `te` is not "two" but a plural marker?} \\
\texttt{{-} Let's look at "a bird chase two stones" = `ketoge te lupohu pepa`.} \\
\texttt{{-} V=`ketoge`. O=`te lupohu`? S=`pepa`?} \\
\texttt{{-} If V{-}O{-}S: O=`te lupohu` (two stones?), S=`pepa` (bird).} \\
\texttt{{-} If `lupo`=stone, then `lupohu` = stone+hu.} \\
\texttt{{-} "two stones" = `te lupohu`.} \\
\texttt{{-} This implies `te` is "two" and `lupohu` is "stones" (plural stone).} \\
\texttt{{-} But `lupohu` = `lupo` (stone) + `hu`.} \\
\texttt{{-} So `hu` marks plural noun?} \\
\texttt{{-} But in "two birds chase a stone", `pepahu` = `pepa` (bird) + `hu`.} \\
\texttt{{-} So `hu` marks plural noun.} \\
\texttt{{-} What about `te`?} \\
\texttt{{-} "two birds" = `pepa te`.} \\
\texttt{{-} "two stones" = `te lupohu`.} \\
\texttt{{-} This is inconsistent. `te` is after `pepa` but before `lupohu`.} \\
\texttt{{-} Maybe `te` is the number word "two", and it comes after the noun?} \\
\texttt{{-} "two birds" = `pepa te`.} \\
\texttt{{-} "two stones" = `lupo te`? But we have `te lupohu`.} \\
\texttt{{-} Maybe the order is Number{-}Noun? `te` `lupohu`.} \\
\texttt{{-} Then "two birds" should be `te pepa`. But we have `pepa te` in the subject slot of `ketoge lupo te pepahu`.} \\
\texttt{{-} Wait, `ketoge lupo te pepahu`.} \\
\texttt{{-} If V{-}O{-}S: V=`ketoge`, O=`lupo` (stone), S=`pepa te` (two birds) + `hu`?} \\
\texttt{{-} If `hu` is plural marker, then `pepa` (bird) + `te` (two) + `hu` (plural) {-}> `pepahu`? No, `te` is in the middle.} \\
\texttt{{-} Maybe `te` is not "two". Maybe `te` is a plural marker for the object?} \\
\texttt{{-} Let's re{-}evaluate "two birds chase a stone" = `ketoge lupo te pepahu`.} \\
\texttt{{-} If V{-}O{-}S: O=`lupo` (stone). S=`pepa` (bird).} \\
\texttt{{-} `lupo te` = "stone two"? No.} \\
\texttt{{-} Maybe the order is V{-}S{-}O?} \\
\texttt{{-} V=`ketoge`. S=`lupo te` (two birds). O=`pepahu` (stone + hu).} \\
\texttt{{-} If S=`lupo te`, then `lupo`=bird.} \\
\texttt{{-} If O=`pepahu`, then `pepa`=stone + `hu`.} \\
\texttt{{-} "a bird chase two stones" = `ketoge te lupohu pepa`.} \\
\texttt{{-} V=`ketoge`. S=`te lupohu`? O=`pepa`?} \\
\texttt{{-} If S=`te lupohu`, then `te` is before noun.} \\
\texttt{{-} If O=`pepa`, then `pepa`=stone.} \\
\texttt{{-} This implies `lupo`=bird, `pepa`=stone.} \\
\texttt{{-} And `te` is "two".} \\
\texttt{{-} "two birds" = `lupo te` (Noun + Number).} \\
\texttt{{-} "two stones" = `te lupohu` (Number + Noun+hu)? Inconsistent.} \\
\texttt{{-} Unless `lupohu` is "stones" and `lupo` is "bird".} \\
\texttt{{-} If `lupo`=bird, `pepa`=stone.} \\
\texttt{{-} "two birds" = `lupo te`.} \\
\texttt{{-} "two stones" = `pepa te`.} \\
\texttt{{-} "a bird chase two stones" = `ketoge te lupohu pepa`.} \\
\texttt{{-} If V{-}S{-}O: S=`te lupohu` (two stones). O=`pepa` (a bird).} \\
\texttt{{-} This implies `pepa`=bird. Contradiction.} \\
\\ 
\texttt{Let's try V{-}O{-}S again with `lupo`=stone, `pepa`=bird.} \\
\texttt{1. "a bird chase a stone" = `ketoge lupo pepa`.} \\
\texttt{   {-} V=`ketoge`. O=`lupo` (stone). S=`pepa` (bird).} \\
\texttt{   {-} Order: V{-}O{-}S.} \\
\texttt{2. "two birds chase a stone" = `ketoge lupo te pepahu`.} \\
\texttt{   {-} V=`ketoge`. O=`lupo` (stone). S=`pepa te` (two birds) + `hu`?} \\
\texttt{   {-} `pepa` (bird) + `te` (two) + `hu` (plural marker?) {-}> `pepahu`.} \\
\texttt{   {-} Maybe `te` and `hu` combine? Or `te` is "two" and `hu` is plural, and they merge?} \\
\texttt{   {-} Or maybe `te` is the plural marker and `hu` is "two"? No, "two" is a number word.} \\
\texttt{   {-} Let's assume `te` = "two".} \\
\texttt{   {-} "two birds" = `pepa te`.} \\
\texttt{   {-} "two stones" = `lupo te`.} \\
\texttt{   {-} In "two birds chase a stone" (V{-}O{-}S): `ketoge` (V) `lupo` (O) `pepa te` (S).} \\
\texttt{   {-} But the output is `ketoge lupo te pepahu`.} \\
\texttt{   {-} This looks like `lupo` (O) `te` (Number?) `pepahu` (S).} \\
\texttt{   {-} Maybe `te` is a plural marker for the Object?} \\
\texttt{   {-} "two birds chase a stone" {-}> Object is "a stone". Why `te`?} \\
\texttt{   {-} Maybe `te` is the number "two" and it modifies the Subject, but the Subject is at the end?} \\
\texttt{   {-} `pepa` (bird) `te` (two) `hu` (plural marker).} \\
\texttt{   {-} "two birds" = `pepa te hu`? But it's `pepahu`.} \\
\texttt{   {-} Maybe `te` is not "two". Maybe `te` is a plural marker for the noun before it?} \\
\texttt{   {-} `lupo te` = "two stones"? No, `lupo` is stone.} \\
\texttt{   {-} Let's look at "a bird chase two stones" = `ketoge te lupohu pepa`.} \\
\texttt{   {-} V=`ketoge`. O=`te lupohu`. S=`pepa`.} \\
\texttt{   {-} If V{-}O{-}S: O=`te lupohu` (two stones). S=`pepa` (bird).} \\
\texttt{   {-} `te` is before `lupohu`. `lupohu` = `lupo` (stone) + `hu`.} \\
\texttt{   {-} So "two stones" = `te` + `lupo` + `hu`.} \\
\texttt{   {-} "two birds" = `pepa` + `te` + `hu`? {-}> `pepahu`.} \\
\texttt{   {-} This fits!} \\
\texttt{   {-} Rule: Number word `te` ("two") comes BEFORE the noun if the noun is plural?} \\
\texttt{   {-} "two stones" = `te` `lupohu`.} \\
\texttt{   {-} "two birds" = `pepahu` (where `te` is missing or merged?).} \\
\texttt{   {-} Wait, `pepahu` = `pepa` + `hu`. Where is `te`?} \\
\texttt{   {-} Maybe `te` is the plural marker?} \\
\texttt{   {-} "two stones" = `te` `lupohu`. `te` is plural marker?} \\
\texttt{   {-} "two birds" = `pepahu`. `hu` is plural marker?} \\
\texttt{   {-} Inconsistent.} \\
\\ 
\texttt{Let's try a different angle.} \\
\texttt{Maybe the language is V{-}S{-}O.} \\
\texttt{1. "a bird chase a stone" = `ketoge lupo pepa`.} \\
\texttt{   {-} S=`lupo` (bird). O=`pepa` (stone).} \\
\texttt{2. "two birds chase a stone" = `ketoge lupo te pepahu`.} \\
\texttt{   {-} S=`lupo te` (two birds). O=`pepahu` (stone + hu).} \\
\texttt{   {-} `hu` marks the object when the subject is plural?} \\
\texttt{3. "a bird chase two stones" = `ketoge te lupohu pepa`.} \\
\texttt{   {-} S=`te lupohu` (two stones). O=`pepa` (bird).} \\
\texttt{   {-} `te` is before the noun. `lupohu` = `lupo` (bird) + `hu`.} \\
\texttt{   {-} If `lupo`=bird, then `lupohu` = bird+hu.} \\
\texttt{   {-} But "two stones" should be stone+plural.} \\
\texttt{   {-} This implies `lupo` = stone.} \\
\texttt{   {-} If `lupo` = stone, then "a bird chase a stone" = `ketoge lupo pepa` {-}> S=`lupo` (stone), O=`pepa` (bird).} \\
\texttt{   {-} "two birds chase a stone" = `ketoge lupo te pepahu`.} \\
\texttt{   {-} S=`lupo te` (two stones). O=`pepahu` (bird + hu).} \\
\texttt{   {-} English: "two birds chase a stone".} \\
\texttt{   {-} So S=`lupo te` must be "two birds".} \\
\texttt{   {-} This implies `lupo` = bird.} \\
\texttt{   {-} Contradiction again.} \\
\\ 
\texttt{Okay, let's look at the words again.} \\
\texttt{`lupo` and `pepa`.} \\
\texttt{Maybe `lupo` = bird, `pepa` = stone.} \\
\texttt{"a bird chase a stone" = `ketoge lupo pepa`. (V S O).} \\
\texttt{"two birds chase a stone" = `ketoge lupo te pepahu`.} \\
\texttt{{-} S=`lupo te` (two birds).} \\
\texttt{{-} O=`pepahu` (stone + hu).} \\
\texttt{"a bird chase two stones" = `ketoge te lupohu pepa`.} \\
\texttt{{-} S=`te lupohu` (two stones).} \\
\texttt{{-} O=`pepa` (bird).} \\
\texttt{{-} This implies `lupohu` = stones.} \\
\texttt{{-} If `lupo` = bird, then `lupohu` = bird + hu.} \\
\texttt{{-} But "two stones" = `te` + `lupohu`.} \\
\texttt{{-} So `lupohu` must be "stones".} \\
\texttt{{-} This implies `lupo` = stone.} \\
\texttt{{-} But `lupo` was "bird" in sentence 1.} \\
\texttt{{-} Is it possible `lupo` means both? No.} \\
\texttt{{-} Is it possible the order is V{-}O{-}S?} \\
\texttt{{-} "a bird chase a stone" = `ketoge lupo pepa`.} \\
\texttt{  {-} V=`ketoge`. O=`lupo` (stone). S=`pepa` (bird).} \\
\texttt{{-} "two birds chase a stone" = `ketoge lupo te pepahu`.} \\
\texttt{  {-} V=`ketoge`. O=`lupo` (stone). S=`pepa te` (two birds) + `hu`?} \\
\texttt{  {-} `pepa` (bird) + `te` (two) + `hu` (plural marker) {-}> `pepahu`.} \\
\texttt{  {-} This works if `te` and `hu` merge or `te` is "two" and `hu` is plural.} \\
\texttt{{-} "a bird chase two stones" = `ketoge te lupohu pepa`.} \\
\texttt{  {-} V=`ketoge`. O=`te lupohu` (two stones). S=`pepa` (bird).} \\
\texttt{  {-} `te` (two) + `lupohu` (stones).} \\
\texttt{  {-} `lupohu` = `lupo` (stone) + `hu` (plural).} \\
\texttt{  {-} This works!} \\
\texttt{{-} "two stones chase a bird" = `ketoge pepa te lupohu`.} \\
\texttt{  {-} V=`ketoge`. O=`pepa te` (two birds). S=`lupohu` (stone + hu).} \\
\texttt{  {-} `pepa` (bird) + `te` (two) {-}> `pepa te`.} \\
\texttt{  {-} `lupohu` = `lupo` (stone) + `hu`.} \\
\texttt{  {-} This works!} \\
\texttt{{-} "two birds chase two stones" = `ketoge te lupohu te pepahu`.} \\
\texttt{  {-} V=`ketoge`. O=`te lupohu` (two stones). S=`pepa te` (two birds) + `hu` {-}> `pepahu`.} \\
\texttt{  {-} Wait, `pepahu` is at the end.} \\
\texttt{  {-} `te lupohu` (two stones) `te pepahu` (two birds + hu).} \\
\texttt{  {-} This fits V{-}O{-}S.} \\
\texttt{{-} "two birds chase the stone" = `ketoge lupo la te pepahu`.} \\
\texttt{  {-} V=`ketoge`. O=`lupo la` (the stone). S=`pepa te` (two birds) + `hu` {-}> `pepahu`.} \\
\texttt{  {-} `lupo la` = "the stone".} \\
\texttt{  {-} `pepahu` = "two birds" (plural subject marker).} \\
\texttt{  {-} This fits!} \\
\\ 
\texttt{Conclusion:} \\
\texttt{{-} Word Order: V{-}O{-}S (Verb{-}Object{-}Subject).} \\
\texttt{{-} Nouns: `lupo` = stone, `pepa` = bird.} \\
\texttt{{-} Number "two": `te`.} \\
\texttt{  {-} Plural Object: `te` + Noun + `hu`? No.} \\
\texttt{  {-} "two stones" = `te lupohu`. `lupohu` = `lupo` + `hu`.} \\
\texttt{  {-} "two birds" = `pepa te`. (In "two birds chase a stone", S=`pepa te` + `hu` {-}> `pepahu`).} \\
\texttt{  {-} Wait, in "two birds chase a stone", S=`pepa te` + `hu` {-}> `pepahu`.} \\
\texttt{  {-} In "two stones chase a bird", O=`pepa te`.} \\
\texttt{  {-} So `te` is "two".} \\
\texttt{  {-} Plural marker `hu` is added to the noun if it is the Subject and plural?} \\
\texttt{  {-} "two birds" (Subject) {-}> `pepa te` + `hu` {-}> `pepahu`.} \\
\texttt{  {-} "two stones" (Object) {-}> `te lupohu`. `lupohu` = `lupo` + `hu`.} \\
\texttt{  {-} So `hu` is added to the noun regardless of position if it's plural?} \\
\texttt{  {-} "two stones" = `te` + `lupo` + `hu`.} \\
\texttt{  {-} "two birds" = `pepa` + `te` + `hu`? But we see `pepahu`.} \\
\texttt{  {-} Maybe `te` and `hu` merge? Or `te` is after the noun for birds?} \\
\texttt{  {-} "two birds" = `pepa te`. With `hu` {-}> `pepahu`.} \\
\texttt{  {-} "two stones" = `te lupohu`.} \\
\texttt{  {-} This is inconsistent.} \\
\texttt{  {-} Let's look at "two birds chase two stones" = `ketoge te lupohu te pepahu`.} \\
\texttt{  {-} O=`te lupohu`. S=`te pepahu`.} \\
\texttt{  {-} So "two birds" = `te pepahu`.} \\
\texttt{  {-} "two stones" = `te lupohu`.} \\
\texttt{  {-} So `te` is "two".} \\
\texttt{  {-} `hu` is plural marker.} \\
\texttt{  {-} For birds: `pepa` + `te` + `hu` {-}> `te` `pepa` `hu`? No, `te` `pepahu`.} \\
\texttt{  {-} For stones: `te` `lupohu`.} \\
\texttt{  {-} Maybe `te` moves?} \\
\texttt{  {-} Or maybe `pepa` is plural by default? No, "a bird" is `pepa`.} \\
\texttt{  {-} Let's assume the rule is:} \\
\texttt{    {-} Plural Noun = Noun + `hu`.} \\
\texttt{    {-} Number "two" = `te`.} \\
\texttt{    {-} If "two" + Noun:} \\
\texttt{      {-} If Noun is stone: `te` + `lupo` + `hu` = `te lupohu`.} \\
\texttt{      {-} If Noun is bird: `pepa` + `te` + `hu` = `te pepahu`?} \\
\texttt{      {-} In "two birds chase two stones", we have `te pepahu`.} \\
\texttt{      {-} So "two birds" = `te pepahu`.} \\
\texttt{      {-} "two stones" = `te lupohu`.} \\
\texttt{      {-} This is consistent! `te` comes before the plural noun phrase.} \\
\texttt{      {-} Plural noun = Noun + `hu`.} \\
\texttt{      {-} "two birds" = `te` + `pepa` + `hu` = `te pepahu`.} \\
\texttt{      {-} "two stones" = `te` + `lupo` + `hu` = `te lupohu`.} \\
\texttt{      {-} But in "two birds chase a stone" = `ketoge lupo te pepahu`.} \\
\texttt{      {-} Here "two birds" is the Subject.} \\
\texttt{      {-} V=`ketoge`. O=`lupo` (stone). S=`te pepahu`?} \\
\texttt{      {-} But the output is `lupo te pepahu`.} \\
\texttt{      {-} `lupo` (stone) `te` (two) `pepahu` (birds).} \\
\texttt{      {-} This suggests `te` is between O and S?} \\
\texttt{      {-} Or maybe `te` is part of the Subject but placed after the Object?} \\
\texttt{      {-} "two birds" = `te pepahu`.} \\
\texttt{      {-} In "two birds chase a stone", the Subject is `te pepahu`.} \\
\texttt{      {-} But the output is `lupo te pepahu`.} \\
\texttt{      {-} Maybe `te` is a plural marker for the Object?} \\
\texttt{      {-} "a stone" = `lupo`. "two stones" = `te lupohu`.} \\
\texttt{      {-} "a bird" = `pepa`. "two birds" = `te pepahu`.} \\
\texttt{      {-} In "two birds chase a stone": V=`ketoge`. O=`lupo`. S=`te pepahu`.} \\
\texttt{      {-} Output: `ketoge lupo te pepahu`.} \\
\texttt{      {-} This matches V{-}O{-}S with S=`te pepahu`.} \\
\texttt{      {-} In "a bird chase two stones": V=`ketoge`. O=`te lupohu`. S=`pepa`.} \\
\texttt{      {-} Output: `ketoge te lupohu pepa`.} \\
            }
        }
    }
    & & \\ \\

    \theutterance \stepcounter{utterance}  
    & & & \multicolumn{2}{p{0.3\linewidth}}{
        \cellcolor[rgb]{0.9,0.9,0.9}{
            \makecell[{{p{\linewidth}}}]{
                \texttt{\tiny{[GM$|$GM]}}
                \texttt{could not find any probe sentences; put them under a 'PROBES:' line, one per line} \\
            }
        }
    }
    & & \\ \\

\end{supertabular}
}

\end{document}
